\chapter{Epilogue: One Assumption, One Principle, One Equation}

Every chapter in this book has followed the same discipline: state what
is proved, state what is conjectured, and do not conflate the two. This
final chapter applies that same discipline to the framework as a whole.

We began with a bold claim: assume nothing, derive everything. The claim
did not survive contact with its own logic. Assuming nothing is itself
an assumption. Any axiomatic system requires at least one starting
point. The question is not whether to make an assumption but which
assumption is minimal.

The answer the framework converges on is this: one assumption, one
parameter, one equation.

\section{The Minimum Assumption}

The single assumption is the existence of an observer.

Not a specific observer. Not a human, not a brain, not a physical
system of any particular kind. The assumption is only that there exists
at least one finite information structure whose state sequence encodes
causally efficacious subjective experience --- the operational
definition established in Chapter~1. We don't argue to know what it is,
just that there is one.

This is not an empty assumption. Chapter~1 proved from four modest
empirical axioms that such a structure must be finite, substrate
independent, and that the universe containing it must be static and
timeless. The existence of the observer forces the existence of a finite
bit count $n$. The universe is the minimal information structure that
contains the observer.

Everything else in the framework is a consequence of this starting
point, combined with one principle: structures are weighted by their
compressibility.

\section{The Organising Equation}

Given an observer $O$, the probability of a history $\gamma$ compatible
with $O$'s existence is:

\begin{equation}
\mathbb{P}(\gamma \mid O)
= \frac{1}{Z_O}\exp\!\bigl(-\mathcal{C}_O[\gamma]\bigr),
\qquad \gamma \in \Gamma_O,
\label{eq:master}
\end{equation}

where $\mathcal{C}_O[\gamma]$ is the total description cost of history
$\gamma$ given observer $O$, $\Gamma_O$ is the set of all histories
compatible with $O$'s existence, and $Z_O$ is the normalisation.

This is Solomonoff induction in path-integral form. The apparent
inverse-temperature parameter has been set to unity, which is not
fine-tuning: rescaling it merely redefines the units of description
cost. Only the ordering of histories by compressibility is physically
meaningful.

The physical laws we observe are the large-deviation minimisers of
$\mathcal{C}_O$:

\[
\gamma^* = \arg\min_{\gamma \in \Gamma_O} \mathcal{C}_O[\gamma].
\]

A universe governed by predictable, differentiable laws is not
mandated from outside. It is selected because ordered configurations
are shorter to describe than chaotic ones. Chaotic universes are not
forbidden; they are exponentially suppressed.

\section{What the Framework Has Established}

Ten chapters of exact results support this picture. We state them
honestly, separating what is proved from what is conjectured.

\subsection*{Exact results}

These results are analytically derived and numerically verified. They
require no free parameters beyond the definition of the density matrix
decomposition.

The diagonal/off-diagonal split $\rho \mapsto \mathrm{diag}(\rho) +
(\rho - \mathrm{diag}(\rho))$ applied to fermionic configurations
gives: Pauli exclusion as zero compression residual; boson amplitude
$\|B(\theta)\| = \sin(2\theta)/\sqrt{2}$ independent of system size;
the Born rule as the interference term of the density matrix; the
virtual propagator lifecycle without field theory.

The same split applied to metric configurations gives: the Ricci/Weyl
decomposition with tracelessness exact; two graviton polarisation modes
as degenerate eigenmodes; graviton amplitude $\|W(\epsilon)\| =
\sin(2\epsilon)/\sqrt{2}$; the Newtonian potential $V(R) = -GMm/R$
with $G = 1/(8\pi)$ in Planck units from flux conservation alone;
Keplerian orbits as minimum-complexity closed worldlines.

The conservation law $\|\mathrm{local}\|^2 + \|\mathrm{non-local}\|^2
= 1$ holds for fermionic and metric configurations alike, as a direct
consequence of $\|\rho\|^2 = 1$ for pure states.

The aspect ratio $\mathcal{A}(n) = 2^n/(n\ln n)$ calibrated against
inflationary expansion in Planck units gives $n \approx 184$, with the
e-fold count and Bekenstein-Hawking entropy recovered as independent
consistency checks.

The Spearman correlation between spectral complexity $C_s$ and
Euclidean action $S_E$ in Wheeler-DeWitt minisuperspace exceeds
$\rho = 0.96$ across 20,000 paths, with the partial correlation
surviving explicit roughness confound control. The Hartle-Hawking
instanton is the joint global minimum of both quantities.

\subsection*{Strongly suggested but not yet derived}

The three-phase expansion profile of standard cosmology --- inflation,
deceleration, re-acceleration --- emerges from the relational scale
factor under lognormal matter abundance without a cosmological constant.
The Hubble tension is a ranging artefact of a system whose
quantisation density has evolved between measurement epochs. The
$4\pi$ factor appears independently in Bekenstein-Hawking recovery, the
Newtonian potential, and the aspect ratio calculation, and in each case
it is the unique conversion between flat counting and spherical
propagation.

\subsection*{Open}

The saddle-point equation

\[
\frac{\partial}{\partial n}
\Bigl[\log Z_O(n) - \log Z(n) - \log n\Bigr] = 0
\]

should select $n^* \approx 184$ from the requirement that the universe
be large enough to contain a self-describing observer yet small enough
to remain compressible. This has not been verified because no tractable
expression for $Z_O(n)$ is currently available. This is the central
open problem of the programme.

The D-$\psi$-G decomposition of the cost functional

\[
\mathcal{C}_O[\gamma] = C_D[n] + C_\psi[\psi_\gamma] + C_G[g_\gamma]
\]

is conjectured, not derived. The three regimes --- discrete at Planck
scale, spectral at quantum scale, geometric at macroscopic scale ---
are motivated by the exact results but have not been shown to follow
from equation~\eqref{eq:master} without additional assumptions.

The continuum limit: recovering the full Einstein field equations from
the discrete toy model of Chapter~9. The derivation of massive bosons,
the gauge structure of the Standard Model, and the quark mass
hierarchy. The electromagnetic force law from the $n=2$, $m=1$ fermion
residual. The proportionality $\|W\| \propto M$ needed to close the
Newtonian potential theorem.

\section{The Shared Formula}

The most striking result across the series is that

\[
\frac{\sin(2\alpha)}{\sqrt{2}}
\]

governs both the boson amplitude (Chapter~7) and the graviton
amplitude (Chapter~9), for entirely different physical degrees of
freedom.

The explanation is algebraic. For any state
$\psi = \cos\alpha\,u + i\sin\alpha\,v$ with $u \perp v$ and $\|u\|
= \|v\| = 1$, the off-diagonal norm of $\rho = |\psi\rangle\langle\psi|$
satisfies

\[
\|\rho - \mathrm{diag}(\rho)\|_F = \frac{|\sin 2\alpha|}{\sqrt{2}}
\]

as a pure algebraic identity, independent of what $u$ and $v$
represent. The formula is not a feature of quantum mechanics or of
gravity. It is a feature of the density matrix decomposition applied
to any normalised superposition of two orthogonal configurations.

Quantum mechanics and general relativity share this formula because
they are both projections of the same compression operation onto
different physical degrees of freedom. This is the conjecture the
framework is built to prove.

\section{What Is Not Claimed}

The framework is not a completed theory. It is a collection of exact
results in discrete toy models, a set of structural correspondences
with known physics, and a conjectured organising equation.

The toy models use one-dimensional chains, scalar conformal metric
factors, and two-frame sequences. Extensions to four continuous
dimensions, the full rank-4 Riemann tensor, and the Fock space of
many-fermion systems remain to be carried out. The framework does not
yet derive the masses of elementary particles, the coupling constants
of the Standard Model, or the value of $\hbar$ in physical units from
$n$ alone.

The claim that the framework is a zero-parameter theory is too strong.
It has one parameter in the sense that matters: the definition of the
observer. Given that definition, $n$ is expected to be fixed by the
saddle-point equation rather than fitted to data. Whether that
expectation survives rigorous calculation is an open question.

What can be said honestly is this: the results accumulated across ten
chapters were not put in by hand. Pauli exclusion, the Born rule, the
Ricci/Weyl split, gravitational wave polarisations, the Newtonian
potential, Keplerian orbits, inflationary expansion, and
Bekenstein-Hawking entropy all follow from the single operation
$B = \rho - \mathrm{diag}(\rho)$ applied to compressed configurations,
combined with the principle that compressible histories dominate the
measure.

Whether the principle alone is sufficient to derive everything
that remains open is the question the research programme exists to
answer.

\section{The Path Forward}

Three problems are foundational. Everything else depends on them.

\textbf{Third}: take the continuum limit of the discrete metric toy
model and recover the Einstein field equations as the stationarity
conditions of $C_G$ over metric configurations. This is the bridge
between the exact toy-model results of Chapter~9 and the full
apparatus of general relativity.

\textbf{Second}: solve the saddle-point equation and recover $n^*
\approx 184$. This requires a tractable model of the observer partition
function $Z_O(n)$. Even a toy model that captures the qualitative
dependence --- too few bits cannot support a self-describing observer,
too many bits drown the ordered histories in a sea of chaos --- would
be a significant step.

\textbf{Third}: derive the D-$\psi$-G decomposition from
equation~\eqref{eq:master}. Show that the boundary structure of any
self-describing observer forces the cost functional to split into
discrete, spectral, and geometric regimes. This would convert the
trinity from a conjecture to a theorem.

If these three problems are solved, the remaining open questions ---
the Standard Model gauge structure, the quark mass hierarchy, the
derivation of $\hbar$ --- become calculations within an established
framework rather than foundational challenges.

\section{A Final Remark}

The project began in 2001 with a simple observation: a computer running
a simulation and the simulation itself are two arrangements of the same
information. Neither is more fundamental. The observer is an
equivalence class of descriptions, not any particular description.

Twenty-five years of development have not changed that starting point.
They have only made it more precise. The observer is the invariant. The
density matrix is the equivalence class representative. The Born rule
is the probability measure on the class. The wavefunction is the codec
that selects which representative the observer experiences.

Quantum mechanics and general relativity are two ways of reading the
same compressed structure. One reads the pixels. The other reads the
geometry they form. The codec is the same in both cases.

That is the claim. The proof is still being written.
